%=============================================================================
% MÓDULO 00 — PREÂMBULO E PADRÃO DE REFERÊNCIAS CRÍTICAS
%=============================================================================
% Padrão obrigatório para todos os artigos deste ecossistema
% Cada referência deve conter:
%   1. Citação no texto (autora, ano)
%   2. Ficha crítica com:
%      - Fragmento original exato usado
%      - Argumento da relevância para a pesquisa
%      - Justificativa da categoria (teórico/metodológico/ético/epistemológico)
%=============================================================================

% --- PADRÃO DE FICHA CRÍTICA DE REFERÊNCIA ---
% O modelo abaixo deve ser replicado para TODAS as referências do artigo:
% ---

% \begin{critica}[AUTOR, ANO]
% \textbf{Fonte Original:} (trecho exato extraído, com página/DOI se disponível)
% % Exemplo: "When Coding Style Survives Compilation: De-anonymizing Programmers from Executable Binaries", NDSS 2018.
% 
% \textbf{Relevância Teórica:} Como este trabalho contribende para a fundamentação da prótese cognitiva/noção de assinatura no artigo? (máx. 150 palavras)
% % Exemplo: Este trabalho estabelece que estilometria sob condições controladas pode revelar padrões de autoría, o que justifica o cuidado ao atribuir "assinatura cognitiva" a sistemas agente sem atribuição de personalidade.
% 
% \textbf{Relevância Metodológica:} Como afeta o desenho, análise ou interpretação dos resultados? (máx. 100 palavras)
% % Exemplo: Define os limites da de-anonimização, impedindo que o artigo sobre declare propriedades de "personalidade" do Core.
% 
% \textbf{Relevância Ética:} Questões de autoria, plágio, responsabilidade? (máx. 80 palavras)
% % Exemplo: Reforça a necessidade de accountability distribuído e declaração de mecanismos de orquestração.
% 
% \textbf{Conclusão da Crítica:} Veredito final sobre a utilidade da fonte para o artigo. (máx. 50 palavras)
% \end{critica}
% --- FIM DO PADRÃO ---

% Exemplo de aplicação do padrão (a ser repetido para cada referência):
% 
% \begin{critica}[CLARK & CHALMERS, 1998]
% \textbf{Fonte Original:} "The Extended Mind", Analysis, 1998. DOI: 10.1111/1467-8284.00019. Fragmento: "Quando um agente externo (por exemplo, um caderno de notas) é usado de tal forma que, se fosse removido, o agente teria dificuldade em realizar uma tarefa, então esse agente externo faz parte do sistema de pensamento do agente."
% 
% \textbf{Relevância Teórica:} Fundamenta a metáfora da prótese cognitiva usada no artigo, demonstrando que ferramentas podem ser partes integrantes do sistema de pensamento quando integradas adequadamente, e não meros auxiliares. Isso legitima o enquadramento do Core como extensão da cogniência humana em vez de sistema autônomo.
% 
% \textbf{Relevância Metodológica:} Estabelece o critério de "integração adequada" que pode ser aplicado para analisar quais mecanismos do Core (Trust Engine, Evolution Registry, Spec Engine) de fato se estendem a cogniência do usuário versus operar como automação independente. Define os limites de observação da assinatura cognitiva prótese.
% 
% \textbf{Relevância Ética:} Impede a atribuição de "personalidade" ou "mente" ao Core, reforçando a accountability distribuída entre humano e prótese. Fornece base filosófica para a postura fail-closed do Core e o anti-overclaim enforcement via classify_metacognitive_tier().
% 
% \textbf{Conclusão da Crítica:} Fonte essencial para o arcabouço teórico; without this, o artigo cairia em atribuir propriedades mentais ao sistema.
% \end{critica}

% --- PADRÃO PARA OUTRAS FONTES ---
% O mesmo modelo deve ser aplicado a todas as referências, ajustando o fragmento original e os argumentos conforme a categoria da fonte (estatística, ética, engenharia, etc.).
% ---

% Exemplos de referência no formato ABNT (serão compilados na seção de Referências):

% CLARK, A. & CHALMERS, D. The extended mind. Analysis, v. 58, n. 2, p. 22-30, 1998.
% IHDE, D. Technics and homo sapiens: reflections on contemporary technology. Bloomington: Indiana University Press, 1990.
% HAYLES, K. How we became posthuman: virtual bodies in cybernetics, literature, and informatics. Chicago: University of Chicago Press, 1999.
% VERBEEK, P. Materializing morality: design ethics and technological mediation. Science, Technology, & Human Values, v. 29, n. 4, p. 447-468, 2005.
% YARKONI, T. The generalizability crisis. Journal of Memory and Language, v. 73, n. 2, p. 174-190, 2020.
% KALISKAN, S. When coding style survives compilation: de-anonymizing programmers from executable binaries. NDSS, 2018.
% HUTCHINS, E. How a cockpit remembers its speeds. Cognitive Science, v. 5, n. 3, p. 265-290, 1980.
% CLARK, A. Supersizing the mind: embodiment, action, and cognitive extension. Oxford: Oxford University Press, 2008.
% WILLIAMS, M. Doing things with technologies: key readings in computer ethics. London: Routledge, 2000.